\chapter{Modelo Físico, Indexação e Desempenho}
\label{cap_modelo_fisico}

A implementação físico-operacional do banco de dados do SIGEA-GTT foi projetada especificamente para o Sistema Gerenciador de Banco de Dados relacional de código aberto de padrão industrial **PostgreSQL 16 LTS** \cite{postgresql2024}. Esta seção aborda as decisões de arquitetura física de dados voltadas para alta disponibilidade, segurança contra ataques de enumeração, integridade transacional e tempos de resposta submétricos.

% ===========================================================================
\section{Estratégia de Chaves Primárias: UUID v4 vs. Sequenciais}
\label{sec_uuid_vs_seq}
% ===========================================================================

Em contraste com abordagens legadas baseadas em inteiros auto-incrementais (\texttt{SERIAL} ou \texttt{BIGSERIAL}), o SIGEA-GTT adota universalmente chaves primárias baseadas em identificadores únicos universais de versão 4 (\textbf{UUID v4} gerados pelo módulo \texttt{uuid-ossp}). A Tabela \ref{tab_comparativo_chaves} resume os fundamentos técnicos dessa decisão:

\begin{table}[!ht]
\centering
\caption{Comparativo entre Chaves Sequenciais e UUID v4}
\label{tab_comparativo_chaves}
\small
\begin{tabularx}{\textwidth}{p{3.2cm}X X}
\toprule
\textbf{Critério de Avaliação} & \textbf{Chaves Inteiras Auto-Incrementais} & \textbf{Chaves UUID v4 Adotadas} \\
\midrule
\textbf{Segurança contra Enumeração} & Vulnerável a ataques \textit{Insecure Direct Object Reference} (IDOR): URLs previsíveis como \texttt{/turmas/12} facilitam raspagem de dados. & Impossível de adivinhar ou enumerar: 128 bits pseudoaleatórios garantem espaço amostral de $3.4 \times 10^{38}$ combinações. \\
\addlinespace
\textbf{Geração Descentralizada} & Centralizada obrigatoriamente no motor do SGBD via \textit{sequences}, gerando contenção em cargas massivas. & Podem ser geradas com segurança criptográfica tanto pelo backend Spring Boot quanto pelo PostgreSQL. \\
\addlinespace
\textbf{Isolamento e LGPD} & Identificadores sequenciais expõem indiretamente a volumetria de discentes e prontuários cadastrados. & Total anonimização referencial, em conformidade com as diretrizes de privacidade por padrão (\textit{privacy by design}). \\
\bottomrule
\end{tabularx}
\fonte{Elaborado pelo autor (2026).}
\end{table}

% ===========================================================================
\section{Estratégia de Indexação e Planos de Execução}
\label{sec_estrategia_indexacao}
% ===========================================================================

Para garantir que o tempo de resposta das consultas analíticas permaneça estritamente inferior a 500 milissegundos (\texttt{RNF03}), definiu-se uma malha de índices B-Tree e GIN detalhada na Tabela \ref{tab_indices_banco}:

\begin{table}[!ht]
\centering
\caption{Relação Canônica de Índices de Desempenho do Banco de Dados}
\label{tab_indices_banco}
\footnotesize
\begin{tabularx}{\textwidth}{p{4.5cm} p{3.9cm} p{1.1cm} X}
\toprule
\textbf{Identificador do Índice} & \textbf{Tabela e Coluna Alvo} & \textbf{Tipo} & \textbf{Finalidade Operacional e Justificativa} \\
\midrule
\texttt{idx\_users\_email} & \texttt{users(email)} & B-Tree & Otimização do fluxo de autenticação e busca por e-mail institucional. \\
\texttt{idx\_users\_role} & \texttt{users(role)} & B-Tree & Aceleração de filtros de perfil nas telas administrativas (\texttt{ADMIN}, \texttt{PROFESSOR}). \\
\texttt{idx\_gtt\_modules\_code} & \texttt{gtt\_modules(code)} & B-Tree & Busca rápida por mnemônico do IHI (\texttt{C}, \texttt{S}, \texttt{M}, \texttt{I}, \texttt{P}, \texttt{E}). \\
\texttt{idx\_gtt\_triggers\_code} & \texttt{gtt\_triggers(code)} & B-Tree & Indexação de busca exata por código alfanumérico do gatilho clínico. \\
\texttt{idx\_gtt\_triggers\_module} & \texttt{gtt\_triggers}\newline\texttt{(module\_id)} & B-Tree & Aceleração da junção referencial entre módulos e seus 53 gatilhos. \\
\texttt{idx\_harm\_sev\_letter} & \texttt{harm\_severities}\newline\texttt{(category\_letter)} & B-Tree & Busca direta pela escala NCC MERP (\texttt{'A'} a \texttt{'I'}). \\
\texttt{idx\_classes\_professor} & \texttt{academic\_classes}\newline\texttt{(professor\_id)} & B-Tree & Filtragem do painel docente com as turmas sob sua titularidade. \\
\texttt{idx\_classes\_period} & \texttt{academic\_classes}\newline\texttt{(academic\_period)} & B-Tree & Filtros de listagem rápida por período letivo semestral. \\
\texttt{idx\_class\_students\_}\newline\texttt{student} & \texttt{class\_students}\newline\texttt{(student\_id)} & B-Tree & Recuperação das turmas em que um estudante está matriculado. \\
\texttt{idx\_activities\_class} & \texttt{activities}\newline\texttt{(class\_id)} & B-Tree & Aceleração de listagem das atividades vinculadas a uma turma. \\
\texttt{idx\_subm\_activity} & \texttt{activity\_submissions}\newline\texttt{(activity\_id)} & B-Tree & Montagem do painel docente de acompanhamento de entregas. \\
\texttt{idx\_subm\_student} & \texttt{activity\_submissions}\newline\texttt{(student\_id)} & B-Tree & Histórico de notas e submissões do estudante logado. \\
\bottomrule
\end{tabularx}
\fonte{Elaborado pelo autor (2026).}
\end{table}

\subsection{Índices GIN para Colunas JSONB}
Para habilitar consultas eficientes no interior de estruturas documentais semiestruturadas, o PostgreSQL disponibiliza os índices invertidos generalizados (\textbf{GIN}). No SIGEA-GTT, esses índices suportam a pesquisa rápida por gatilhos específicos acionados em auditorias anteriores por meio do operador de contenção documental \texttt{@>}:

\begin{lstlisting}[language=SQL, caption={Criação de Índice GIN para Gatilhos Auditados}]
CREATE INDEX idx_submissions_triggers_gin 
ON activity_submissions USING GIN (identified_triggers jsonb_path_ops);
\end{lstlisting}

Essa estrutura permite que a identificação de quais auditorias discentes detectaram o gatilho \texttt{"M03"} ocorra em tempo logarítmico, sem varredura sequencial da tabela (\textit{Seq Scan}).

% ===========================================================================
\section{Políticas de Integridade Referencial (RESTRICT vs. CASCADE)}
\label{sec_politicas_integridade}
% ===========================================================================

A integridade referencial foi calibrada com precisão entre duas políticas fundamentais:

\begin{itemize}
    \item \textbf{Política \texttt{ON DELETE RESTRICT} (Preservação Histórica e Auditoria)}:
    \begin{itemize}
        \item Aplicada a \texttt{users(id)} em todas as referências (\texttt{professor\_id}, \texttt{student\_id}): impede a exclusão acidental de docentes ou discentes que possuam histórico escolar, turmas criadas ou submissões avaliadas registradas no sistema;
        \item Aplicada a \texttt{gtt\_modules(id)} em \texttt{gtt\_triggers}: garante a inviolabilidade da taxonomia internacional do IHI, impedindo a supressão de módulos que possuam gatilhos vinculados.
    \end{itemize}

    \item \textbf{Política \texttt{ON DELETE CASCADE} (Composição Estrutural e Ciclo de Vida)}:
    \begin{itemize}
        \item Aplicada ao vínculo de \texttt{academic\_classes} com \texttt{class\_students} e \texttt{activities}: caso uma turma em fase de rascunho seja descartada antes do início das aulas, suas dependências diretas de composição são limpas atomicamente;
        \item Aplicada ao vínculo de \texttt{activities} com \texttt{activity\_submissions}: se uma atividade for cancelada pela coordenação antes da homologação de notas, os registros vinculados são removidos em cascata de forma transacional.
    \end{itemize}
\end{itemize}
